% Portable starter following the public TEMat article structure.
\documentclass[11pt,a4paper]{article}
\usepackage[utf8]{inputenc}
\usepackage[T1]{fontenc}
\usepackage[spanish,english]{babel}
\usepackage[margin=2.7cm]{geometry}
\usepackage{amsmath,amssymb,amsthm,booktabs,enumitem,xcolor,hyperref}
\hypersetup{colorlinks=true,linkcolor=blue!50!black,urlcolor=blue!50!black}
\newtheorem{proposition}{Proposición}
\newtheorem{example}{Ejemplo}
\title{Un modelo discreto para estudiar redes de transporte}
\author{Lucía Fernández\thanks{Universidad de Ejemplo}\and André Martins}
\date{TEMat -- artículo de muestra}

\begin{document}
\maketitle
\begin{abstract}
Presentamos un modelo discreto para analizar cómo una perturbación se propaga
en una red pequeña. El modelo separa la topología de la regla de actualización,
lo que permite comparar configuraciones sin cambiar el procedimiento de
cálculo. Se incluyen una demostración, un ejemplo numérico y observaciones
para futuras extensiones.
\end{abstract}
\noindent\textbf{Palabras clave:} grafos; sistemas discretos; redes; modelización

\selectlanguage{english}
\begin{abstract}
We present a discrete model for studying how a perturbation propagates through
a small network. The model separates topology from the update rule, allowing
configurations to be compared without changing the calculation procedure.
\end{abstract}
\selectlanguage{spanish}

\section{Introducción}
Los modelos sencillos permiten aislar una hipótesis antes de incorporar
detalles de una aplicación. Sea $G=(V,E)$ un grafo finito y sea $x_t(v)$ el
estado del vértice $v$ en el instante $t$.

\section{Modelo}
Definimos la actualización por
\begin{equation}
 x_{t+1}(v)=\frac{1}{\deg(v)}\sum_{u\sim v}x_t(u),
 \label{eq:update}
\end{equation}
para vértices con grado positivo. La ecuación~\ref{eq:update} representa un
promedio local; en una aplicación real debe justificarse la elección de los
pesos.

\begin{proposition}
Si $G$ es regular y conexo, el promedio global de los estados se conserva tras
una actualización.
\end{proposition}
\begin{proof}
En un grafo regular cada vértice recibe el mismo peso de sus vecinos. Al sumar
la ecuación~\ref{eq:update} sobre todos los vértices, cada estado aparece el
mismo número de veces en numerador y denominador. Por tanto, la suma y el
promedio se conservan.
\end{proof}

\section{Ejemplo}
\begin{example}
Para un ciclo de cuatro vértices, el estado inicial $(1,0,0,0)$ produce
$(0,1/2,0,1/2)$ después de una actualización y $(1/2,0,1/2,0)$ después de
dos. La masa total permanece igual a uno.
\end{example}

\begin{table}[ht]
\centering
\caption{Iteraciones del ejemplo.}
\begin{tabular}{crrrr}
\toprule
Iteración & $x_1$ & $x_2$ & $x_3$ & $x_4$\\
\midrule
0 & 1 & 0 & 0 & 0\\
1 & 0 & 0.5 & 0 & 0.5\\
2 & 0.5 & 0 & 0.5 & 0\\
\bottomrule
\end{tabular}
\end{table}

\section{Conclusiones}
La separación entre grafo y regla de actualización facilita experimentar con
redes alternativas. El texto de ejemplo debe sustituirse por el resultado y
las referencias del trabajo que se desea enviar.

\appendix
\section{Notas de implementación}
Para reproducir los cálculos, registre la representación del grafo, la
precisión numérica y la condición de parada. Mantenga los datos de prueba junto
con el código que genera las tablas.

\end{document}
